\documentclass[11pt,a4paper]{article}
\usepackage[utf8]{inputenc}
\usepackage[T1]{fontenc}
\usepackage[portuguese]{babel}
\usepackage{graphicx}
\usepackage{enumitem}
\usepackage{multicol}
\usepackage{booktabs}
\usepackage{xcolor}
\usepackage{titlesec}
\usepackage{fancyhdr}
\usepackage{amsmath}
\usepackage{geometry}
\usepackage{caption}
\usepackage{microtype}
\usepackage{underscore}


% --- Definição de Cores ---
\definecolor{primaryColor}{RGB}{26, 54, 93}
\definecolor{secondaryColor}{RGB}{43, 108, 176}
\definecolor{textColor}{RGB}{45, 55, 72}
\definecolor{lightGray}{RGB}{247, 250, 252}
\usepackage[colorlinks=true, linkcolor=primaryColor, urlcolor=secondaryColor]{hyperref}

% --- Configuração das Margens ---
\geometry{
    a4paper,
    top=30mm,
    bottom=30mm,
    left=25mm,
    right=25mm,
    headheight=15pt,
    footskip=15mm
}

\color{textColor}
\renewcommand{\baselinestretch}{1.15}

% --- Configuração de Cabeçalhos e Rodapés ---
\pagestyle{fancy}
\fancyhf{}
\fancyhead[L]{\small\color{gray}Mini-Project 1 -- MDRS}
\fancyfoot[C]{\small\color{textColor}\thepage}
\renewcommand{\headrulewidth}{0.5pt}
\renewcommand{\headrule}{\hbox to\headwidth{\color{gray!50}\leaders\hrule height \headrulewidth\hfill}}

% --- Estilização das Secções ---
\titleformat{\section}
  {\Large\bfseries\color{primaryColor}}
  {\thesection}{1em}{}[{\color{secondaryColor!40}\titlerule[1.5pt]}]

\titleformat{\subsection}
  {\large\bfseries\color{secondaryColor}}
  {\thesubsection}{1em}{}

\captionsetup{
    labelfont={bf,color=primaryColor},
    font=small,
    labelsep=quad,
    margin=10pt
}

\setlist[itemize]{leftmargin=*, topsep=4pt, itemsep=3pt, parsep=0pt, partopsep=0pt}

\usepackage{listings}
\definecolor{codegray}{rgb}{0.95,0.95,0.95}

\lstset{
    backgroundcolor=\color{codegray},
    basicstyle=\ttfamily\small,
    frame=single,
    breaklines=true,
    keywordstyle=\color{blue},
    commentstyle=\color{green!50!black},
    stringstyle=\color{red},
    showstringspaces=false,
    tabsize=4
}

\renewcommand{\lstlistingname}{Listagem}

\begin{document}

% ==============================================================================
% PÁGINA DE ROSTO
% ==============================================================================
\begin{titlepage}
    \centering

    \vspace*{1cm}

    {\color{secondaryColor}\rule{\linewidth}{2pt}\par}
    \vspace{0.6cm}

    {\huge \bfseries \color{primaryColor}
    Mini-Project 1:\\
    Performance Evaluation of Point-to-Point Links\par}

    \vspace{0.4cm}

    {\large \color{secondaryColor}
    Supporting Multiple Packet Flows\par}

    \vspace{0.2cm}

    {\color{secondaryColor}\rule{\linewidth}{2pt}\par}

    \vspace{2cm}

    {\Large \textsc{Modelação e Desempenho de Redes e Serviços}\par}

    \vspace{0.3cm}
    {\large Ano letivo 2026/2027\par}

    \vfill

    \begin{tabular}{rlll}
        \textbf{Grupo} & Ricardo Carmo   & (112657) &   50 \%\\
                       & Lázaro Sá       & (115884) &   50 \%\\
                       
    \end{tabular}

    \vspace{0.8cm}

    \vfill

    {\large \color{textColor!70} \today\par}

    \vspace*{1cm}
\end{titlepage}

\clearpage
\tableofcontents
\clearpage

% ==============================================================================
% CONTEÚDO DO RELATÓRIO
% ==============================================================================

\section{Tarefa 1.a -- Simulador para três fluxos}

O simulador das aulas (\texttt{Sim2Flows}, Guião 2) já era genérico em tudo
exceto no número de fluxos: a lista de eventos, a fila em bytes, o teste de
admissão e as fórmulas dos parâmetros de desempenho não dependem do número de
fluxos. A nova versão \texttt{Sim3FlowsT} mantém essa estrutura e generaliza-a
para os três fluxos do enunciado; por isso, apenas se apresentam os trechos
novos ou editados.



\subsection*{Alterações introduzidas e sua justificação}

\begin{itemize}
    \item \textbf{Contadores estatísticos de 2 para 3 posições.} Todos os
    contadores do simulador base (\texttt{TOTALPACKETS}, \texttt{LOSTPACKETS},
    \texttt{TRANSPACKETS}, \texttt{TRANSBYTES}, \texttt{DELAYS},
    \texttt{MAXDELAY}) passaram de \lstinline|zeros(1,2)| para
    \lstinline|zeros(1,3)|. Como as estatísticas são mantidas por fluxo e
    indexadas pelo identificador do fluxo, sem esta alteração uma chegada do
    fluxo~3 originaria um erro de indexação. O mesmo se aplica às saídas
    \texttt{PL}, \texttt{APD}, \texttt{MPD} e \texttt{TT}, que passam a ser
    vetores de três elementos.

    \begin{lstlisting}[basicstyle=\ttfamily\scriptsize]
TOTALPACKETS= zeros(1,3);     % No. of packets of each flow arrived to the system
LOSTPACKETS= zeros(1,3);      % No. of packets of each flow lost due to buffer fullness
TRANSPACKETS= zeros(1,3);     % No. of transmitted packets of each flow
TRANSBYTES= zeros(1,3);       % Sum of the Bytes of all transmitted packets of each flow
DELAYS= zeros(1,3);           % Sum of the delays of all transmitted packets of each flow
MAXDELAY= zeros(1,3);         % Maximum delay among all transmitted packets of each flow
    \end{lstlisting}

    \item \textbf{Chegada inicial do fluxo~3.} Foi acrescentado um terceiro
    evento de chegada inicial,
    \lstinline|tmp = Clock + exprnd(1/lambda(3))|, seguindo exatamente o
    padrão das duas primeiras. Cada fluxo só existe no simulador se tiver uma
    primeira chegada agendada; sem esta linha o fluxo~3 nunca geraria tráfego,
    porque todas as chegadas seguintes são geradas a partir de uma chegada
    anterior do mesmo fluxo.

    \begin{lstlisting}[basicstyle=\ttfamily\scriptsize]
tmp= Clock + exprnd(1/lambda(2));
EventList = [EventList; ARRIVAL, tmp, GenerateSize2(), 2, tmp];
% adicionar o terceiro fluxo
tmp = Clock + exprnd(1/lambda(3));
EventList = [EventList; ARRIVAL, tmp, GenerateSize3(), 3, tmp];
    \end{lstlisting}

    \item \textbf{Selecção do gerador do próximo pacote.} No simulador base a
    escolha era feita com \lstinline|if Flow==1 ... else ...|, o que atribuiria
    ao fluxo~2 o gerador do fluxo~3 (e vice-versa). Foi substituída por um
    \lstinline|switch Flow| com um \lstinline|case| por fluxo, garantindo que
    cada fluxo usa sempre a sua distribuição de tamanhos.

    \begin{lstlisting}[basicstyle=\ttfamily\scriptsize]
            % agora temos um caso para cada fluxo
            switch Flow
                case 1
                    nextPacketSize= GenerateSize1();
                case 2
                    nextPacketSize= GenerateSize2();
                case 3
                    nextPacketSize= GenerateSize3();
            end
    \end{lstlisting}

    \item \textbf{Distribuições dos tamanhos dos pacotes.} Os geradores do
    Guião 2 não correspondem ao enunciado: o fluxo~1 usava a distribuição
    discreta do Guião 1 (64, 110 e 1518~B com probabilidades 0,19, 0,23 e
    0,17, e os restantes valores uniformes entre 65 e 1517~B) e o fluxo~2
    usava 110--130~B. Como o enunciado especifica uma distribuição uniforme
    (valores inteiros com igual probabilidade) em 100--1500~B, 100--150~B e
    500--1000~B, os três geradores passaram a \lstinline|randi([a b])|, que
    devolve um inteiro uniforme no intervalo fechado indicado, e foi criado o
    \lstinline|GenerateSize3|.

    \begin{lstlisting}[basicstyle=\ttfamily\scriptsize]
function out= GenerateSize1()
    out= randi([100 1500]);   % Flow 1: uniforme entre 100 e 1500 bytes
end
function out= GenerateSize2()
    out= randi([100 150]);    % Flow 2: uniforme entre 100 e 150 bytes
end
function out= GenerateSize3()
    out= randi([500 1000]);   % Flow 3: uniforme entre 500 e 1000 bytes
end
    \end{lstlisting}

    \item \textbf{Resto do simulador inalterado.} A lista de eventos mantém as
    cinco colunas (tipo, instante, tamanho, fluxo e instante de chegada), a
    fila continua a ser uma única fila FIFO contabilizada em bytes, o teste de
    admissão continua a ser
    \lstinline|QUEUEOCCUPATION + PacketSize <= f| e as fórmulas de
    \texttt{PL}, \texttt{APD}, \texttt{MPD} e \texttt{TT} não mudaram -- apenas
    passaram a operar sobre vetores de três posições. O identificador do fluxo
    já era guardado na fila e nos eventos, pelo que atribuir as estatísticas ao
    fluxo correto não exigiu qualquer alteração estrutural.
\end{itemize}

O critério de paragem continua a ser \lstinline|sum(TRANSPACKETS)<P|, isto é,
$P$ pacotes transmitidos no total dos três fluxos. Por consequência, o número
de pacotes simulados por fluxo é proporcional à respectiva taxa de chegada.

\newpage

\section{Tarefa 1.b -- Resultados de simulação para o caso C}

Para o caso C ($\lambda_1=\lambda_2=800$ pps, $\lambda_3=600$ pps), com
$f=1\,000\,000$~B e $P=10^5$, obtiveram-se os resultados da
Tabela~\ref{tab:1b}.

\begin{table}[htbp]
    \centering
    \caption{Resultados de simulação para o caso C (50 réplicas, IC de 90\%).}
    \label{tab:1b}
    \begin{tabular}{lccc}
        \toprule
        Parâmetro & Fluxo 1 & Fluxo 2 & Fluxo 3 \\
        \midrule
        PL (\%)   & $0 \pm 0$                & $0 \pm 0$                & $0 \pm 0$ \\
        APD (ms)  & $7.33167 \pm 0.161631$   & $6.79454 \pm 0.159963$   & $7.29599 \pm 0.162140$ \\
        MPD (ms)  & $44.9424 \pm 2.20024$    & $44.5527 \pm 2.21851$    & $44.7092 \pm 2.19380$ \\
        TT (Mbps) & $5.11605 \pm 0.005824$   & $0.799148 \pm 0.001032$  & $3.59820 \pm 0.004634$ \\
        \bottomrule
    \end{tabular}
\end{table}

\subsection*{Análise dos resultados}

\begin{itemize}
    \item \textbf{Não há perdas.} Com $E[L]=800$, $125$ e $750$~B, a carga
    oferecida é de $5.12+0.80+3.60=9.52$~Mbps, ou seja
    $\rho=0.952$: a ligação está muito carregada, mas é estável. O maior atraso
    observado ($\approx 45$~ms) corresponde apenas a
    $45\times10^{-3}\times1.25\times10^{6}\approx 56$~kB de trabalho
    transmitido, uma ordem de grandeza inferior à capacidade da fila de
    1~MB. Esta comparação é apenas indicativa, pois o atraso máximo não
    determina diretamente a ocupação instantânea da fila. A ausência de
    perdas é confirmada diretamente pelos valores de PL, que foram nulos nas
    50 réplicas; o valor $0\pm0$ não resulta de falta de precisão da
    estimativa.

    \item \textbf{O atraso médio é dominado pela espera comum na fila.} Como a
    fila é FIFO e única, todos os pacotes sofrem a mesma espera média; o que
    distingue os fluxos é o tempo de transmissão do próprio pacote,
    $E[S_i]=0.64$, $0.10$ e $0.60$~ms. Ao subtrair estes valores aos atrasos
    simulados obtém-se $6.69$, $6.69$ e $6.70$~ms para os três fluxos, o que
    confirma esta interpretação. O fluxo~2 é o que apresenta menor atraso por
    ter os pacotes mais pequenos (100--150~B).

    \item \textbf{Os atrasos máximos são praticamente iguais.} A diferença entre
    os três fluxos (44.55 a 44.94~ms) é muito inferior à semi-amplitude dos IC
    ($\pm2.2$~ms), logo não é estatisticamente significativa. O máximo é
    determinado pela pior excursão de fila de cada réplica, situação que é
    comum a todos os fluxos por partilharem a mesma fila.

    \item \textbf{O throughput é proporcional a $\lambda_i E[L_i]$.} Não havendo
    perdas, $TT_i=\lambda_i E[L_i]\times8$, o que dá 5.12, 0.80 e 3.60~Mbps: o
    fluxo~1 transporta mais bits por segundo não por ter mais pacotes (tem os
    mesmos 800 pps do fluxo~2), mas por ter pacotes em média seis vezes
    maiores. A soma dos valores de throughput (9.513~Mbps) coincide com a carga oferecida, e
    o IC é muito estreito porque se trata de uma média de bytes ao longo de
    vários segundos de simulação.
\end{itemize}

Em suma, em FIFO, os três fluxos são tratados da mesma forma: o
único fator que os distingue é o tamanho médio dos seus pacotes.

\newpage

\section{Tarefa 1.c -- Modelo M/G/1 para o caso C}

Os três fluxos são multiplexados estatisticamente na mesma ligação, pelo que o
sistema é modelado como uma única fila M/G/1 cujo tempo de serviço é uma
mistura das distribuições dos três fluxos, ponderada pelas taxas de chegada
(Listagem~\ref{lst:1c}).

\begin{lstlisting}[caption={Determinação dos valores teóricos do caso C (M/G/1).},
                   label={lst:1c},
                   basicstyle=\ttfamily\scriptsize]
lambda = [800 800 600];        % packet arrival rates (packets/s)
C = 10e6;                      % link capacity (bit/s)

packetSizes = {100:1500; 100:150; 500:1000};   % uniform discrete sizes (Bytes)

nFlows = numel(lambda);
lambdaTotal = sum(lambda);

% Mean packet size, mean and second moment of the service time of each flow
for flow = 1:nFlows
    sizes = packetSizes{flow};
    serviceTimes = 8 * sizes / C;              % bytes -> bits, C in bit/s
    meanPacketSize(flow) = mean(sizes);
    meanServiceTime(flow) = mean(serviceTimes);
    secondServiceMoment(flow) = mean(serviceTimes.^2);
end

% Service distribution of the aggregate flow (weighted by lambda)
meanAggregateServiceTime = sum(lambda .* meanServiceTime) / lambdaTotal;
secondAggregateServiceMoment = sum(lambda .* secondServiceMoment) / lambdaTotal;

rho = lambdaTotal * meanAggregateServiceTime;  % traffic intensity
if rho >= 1
    error('The M/G/1 system is unstable: rho must be less than 1.');
end

% Pollaczek-Khinchine mean waiting time (common to all flows, FIFO)
averageQueueingDelay = lambdaTotal * secondAggregateServiceMoment / ...
    (2 * (1 - rho));

% Per-flow results: queueing delay + own mean transmission time
averagePacketDelayMs = 1000 * (averageQueueingDelay + meanServiceTime);
throughputMbps = lambda .* meanPacketSize * 8 / 1e6;
\end{lstlisting}

\subsection*{Explicação do código}

\begin{itemize}
    \item \textbf{Momentos do serviço agregado.} Uma chegada pertence ao fluxo
    $i$ com probabilidade $\lambda_i/\Lambda$, com $\Lambda=\sum_i\lambda_i$,
    pelo que
    $E[S]=\sum_i\lambda_i E[S_i]/\Lambda$ e
    $E[S^2]=\sum_i\lambda_i E[S_i^2]/\Lambda$. Os momentos de cada fluxo são
    calculados a partir de todos os valores possíveis do tamanho do pacote
    (distribuição uniforme discreta), e não apenas a partir da média, porque a
    variância do serviço entra na fórmula do tempo de espera. Obtém-se
    $E[S]=0.4327$~ms, $E[S^2]=2.9253\times10^{-7}$~s$^2$ e $\rho=0.952$,
    coerente com a carga oferecida de 9.52~Mbps.

    \item \textbf{Fórmula de Pollaczek--Khinchine.} Para M/G/1,
    $W_q=\Lambda E[S^2]/\left[2(1-\rho)\right]=6.704$~ms. Esta espera é comum
    aos três fluxos, porque a fila é única e FIFO; só depende da carga global.

    \item \textbf{Atraso e throughput por fluxo.} O atraso médio de cada fluxo é
    $W_q+E[S_i]$, somando ao tempo de espera o tempo de transmissão do próprio
    pacote, que depende do tamanho médio do fluxo. O throughput é
    $TT_i=\lambda_iE[L_i]\times8$, uma vez que o modelo admite fila infinita e,
    portanto, ausência de perdas.
\end{itemize}

\subsection*{Resultados e comparação com a simulação}

\begin{table}[htbp]
    \centering
    \caption{Valores teóricos (M/G/1) e simulados para o caso C.}
    \label{tab:1c}
    \begin{tabular}{cccccc}
        \toprule
        Fluxo & $E[L]$ (B) & TT teórico (Mbps) & TT simulado (Mbps) & APD teórico (ms) & APD simulado (ms) \\
        \midrule
        1 & 800 & 5.1200 & $5.1161 \pm 0.0058$ & 7.3438 & $7.3317 \pm 0.1616$ \\
        2 & 125 & 0.8000 & $0.7991 \pm 0.0010$ & 6.8038 & $6.7945 \pm 0.1600$ \\
        3 & 750 & 3.6000 & $3.5982 \pm 0.0046$ & 7.3038 & $7.2960 \pm 0.1621$ \\
        \bottomrule
    \end{tabular}
\end{table}

\begin{table}[htbp]
    \centering
    \caption{Erro relativo do modelo face à simulação.}
    \label{tab:1c-err}
    \begin{tabular}{lccc}
        \toprule
        & Fluxo 1 & Fluxo 2 & Fluxo 3 \\
        \midrule
        Erro relativo do APD & 0.17\% & 0.14\% & 0.11\% \\
        Erro relativo do TT  & 0.08\% & 0.11\% & 0.05\% \\
        \bottomrule
    \end{tabular}
\end{table}

\subsection*{Conclusões}

\begin{itemize}
    \item O modelo M/G/1 reproduz a simulação com erros relativos inferiores a
    0.2\%, bastante abaixo da semi-amplitude dos IC de 90\% ($\pm0.16$~ms em
    atraso e $\pm0.006$~Mbps em throughput). Ou seja, as diferenças observadas não
    são distinguíveis da variabilidade estatística da simulação, o que valida
    quer o simulador quer o modelo.
    \item A estrutura dos resultados é explicada pelo modelo: a espera em fila
    é a mesma para os três fluxos e é a parcela dominante (6.70~ms), enquanto a
    diferença entre fluxos vem apenas do tempo de transmissão médio. Isto
    confirma que, em FIFO, os fluxos não se prejudicam entre si de forma
    seletiva -- todos são penalizados pela carga global.
    \item As pequenas diferenças residuais resultam de a simulação ter um
    horizonte finito ($P=10^5$ pacotes, cerca de 45~s) e de o modelo descrever
    o regime estacionário.
    \item O modelo tem limites que importam para as tarefas seguintes: admite
    fila infinita (não prevê perdas nem o efeito de $f$) e só fornece valores
    médios, pelo que não permite estimar o atraso máximo (MPD). Neste caso,
    como não houve perdas, a hipótese de fila infinita não introduz erro.
\end{itemize}

\newpage

\section{Tarefa 1.d -- Atraso médio e máximo do fluxo 3 (casos A, B e C)}

O script executa \texttt{Sim3FlowsT} 50 vezes para cada um dos casos A, B e C,
com $f=1\,000\,000$~B e $P=10^5$, e apresenta as estimativas do atraso médio e
do atraso máximo do fluxo~3 em gráficos de barras com as barras de erro dos IC
de 90\% (Figuras~\ref{fig:1d-apd} e~\ref{fig:1d-mpd}).

\begin{figure}[htbp]
    \centering
    \IfFileExists{Imagens/Figure_1.png}%
      {\includegraphics[width=0.75\linewidth]{Imagens/Figure_1.png}}%
      {\fbox{\parbox[c][3.5cm][c]{0.7\linewidth}{\centering\small
        Atraso médio do fluxo~3 (adicionar o gráfico exportado do MATLAB em
        \texttt{Imagens/Figure_1.png})}}}
    \caption{Atraso médio do fluxo 3 nos casos A, B e C (estimativa e IC de 90\%).}
    \label{fig:1d-apd}
\end{figure}

\begin{figure}[htbp]
    \centering
    \IfFileExists{Imagens/Figure_2.png}%
      {\includegraphics[width=0.75\linewidth]{Imagens/Figure_2.png}}%
      {\fbox{\parbox[c][3.5cm][c]{0.7\linewidth}{\centering\small
        Atraso máximo do fluxo~3 (adicionar o gráfico exportado do MATLAB em
        \texttt{Imagens/Figure_2.png})}}}
    \caption{Atraso máximo do fluxo 3 nos casos A, B e C (estimativa e IC de 90\%).}
    \label{fig:1d-mpd}
\end{figure}

\subsection*{Justificação dos resultados}

Os três casos diferem apenas na carga: $\rho=0.656$ (A), $0.804$ (B) e
$0.952$ (C). A utilização aparece no denominador $1-\rho$ da fórmula de
Pollaczek--Khinchine, que toma os valores $0.344$, $0.196$ e $0.048$; o tempo
de espera cresce, por isso, muito mais depressa do que a taxa de chegada.
Usando o modelo da Tarefa~1.c, $W_q=0.63$, $1.37$ e $6.70$~ms para A, B e C,
um fator de cerca de 10.6 entre o caso A e o caso C, embora a taxa de
chegada total suba apenas de 1400 para 2200 pps (fator 1.57). Somando o tempo
de transmissão médio do fluxo~3 ($0.60$~ms) obtêm-se atrasos médios teóricos
de 1.23, 1.97 e 7.30~ms; o valor simulado para o caso C,
$7.296\pm0.162$~ms, confirma este comportamento. A passagem de A para B é
moderada, mas de B para C é abrupta: é a zona em que a ligação deixa de ter
folga e a fila passa a acumular pacotes.

O atraso máximo acompanha esta tendência, mas é uma estatística de extremo
(máximo de $10^5$ pacotes por réplica) e por isso muito mais dispersa, como
mostram as barras de erro -- no caso C o IC é de $\pm2.19$~ms, cerca de 5\% do
valor estimado, contra cerca de 2\% no atraso médio. Um único período de
ocupação invulgarmente longo numa réplica altera o máximo, o que explica a
maior variabilidade e o facto de o erro crescer com a carga.

\subsection*{Conclusões}

\begin{itemize}
    \item O atraso do fluxo~3 cresce de forma claramente não linear com a carga
    porque a espera em fila é proporcional a $1/(1-\rho)$; pequenos aumentos de
    tráfego na vizinhança da saturação produzem grandes degradações.
    \item O atraso médio é uma métrica robusta (IC estreitos) e útil para
    comparar cenários, enquanto o atraso máximo, sendo um extremo, exige
    bastantes réplicas para ser estimado com precisão aceitável.
    \item Os casos A, B e C não apresentam perdas: o caso C, o mais carregado,
    já não perde pacotes, pelo que os casos com menor carga também não os
    perdem, e a fila de 1~MB nunca é limitativa nos três cenários.
\end{itemize}

\newpage



\section{Tarefa 2.a -- Simulador para três fluxos com prioridades}

Na Tarefa~2 foi desenvolvida a função \texttt{Sim3FlowsPrioT}, a partir do
simulador \texttt{Sim2FlowsPrio} utilizado no Guião~2. A nova versão passou
a trabalhar com três fluxos e duas filas: a fila~1 é partilhada pelos
fluxos~1 e~2, que têm a mesma prioridade, e a fila~2 contém apenas os pacotes
do fluxo~3. Assim, a extensão não corresponde a duas cópias independentes do
modelo original, mas à junção das duas filas de menor prioridade num único
buffer partilhado:

\begin{itemize}
    \item o fluxo~3 tem a prioridade mais elevada;
    \item os pacotes dos fluxos~1 e~2 ficam na fila~1;
    \item os pacotes do fluxo~3 ficam na fila~2;
    \item a fila~2 é escolhida antes da fila~1 quando \texttt{prio == 1};
    \item dentro da fila dos fluxos~1 e~2, os pacotes são escolhidos pela
    ordem de chegada;
    \item um pacote que já começou a ser transmitido termina antes de se
    escolher outro pacote.
\end{itemize}

O identificador do fluxo continua guardado em cada entrada da fila, o que
permite transmitir pela ordem FIFO e calcular separadamente as perdas, os
atrasos e os \textit{throughputs} dos três fluxos.

\subsection*{Alterações introduzidas}

Relativamente à \texttt{Sim2FlowsPrio}, os contadores estatísticos, os
contadores de chegadas e partidas e a inicialização do evento de chegada do
terceiro fluxo foram adaptados tal como descrito na Tarefa~1 para o acréscimo
de um fluxo.

A diferença específica desta versão é que os fluxos~1 e~2 passam a partilhar
a fila~1, enquanto o fluxo~3 utiliza a fila~2. Por isso, à chegada de um
pacote é atribuído o índice da fila \texttt{Q}, em vez de usar diretamente o
índice \texttt{Flow}:

\begin{lstlisting}[basicstyle=\ttfamily\scriptsize,
                   caption={Encaminhamento dos três fluxos pelas duas filas.},
                   label={lst:2a-routing}]
if Flow == 1 || Flow == 2
    Q = 1;
else
    Q = 2;
end

if QUEUEOCCUPATION(Q) + PacketSize <= f(Q)
    QUEUE{Q} = [QUEUE{Q}; PacketSize, Clock, Flow];
    QUEUEOCCUPATION(Q) = QUEUEOCCUPATION(Q) + PacketSize;
else
    LOSTPACKETS(Flow) = LOSTPACKETS(Flow) + 1;
end
\end{lstlisting}

\begin{lstlisting}[basicstyle=\ttfamily\scriptsize,
                   caption={Seleção da fila no instante da partida.},
                   label={lst:2a-departure}]

QFlow = 0;
if prio == 1
    if QUEUEOCCUPATION(2) > 0
        QFlow = 2;
    elseif QUEUEOCCUPATION(1) > 0
        QFlow = 1;
    else
        STATE = 0;
    end
else
    if QUEUEOCCUPATION(1) > 0
        QFlow = 1;
    elseif QUEUEOCCUPATION(2) > 0
        QFlow = 2;
    else
        STATE = 0;
    end
end

if QFlow > 0
    QSize = QUEUE{QFlow}(1,1);
    QInstant = QUEUE{QFlow}(1,2);
    PacketFlow = QUEUE{QFlow}(1,3);
    EventList = [EventList; DEPARTURE, ...
        Clock + 8*QSize/(C*1e6), QSize, PacketFlow, QInstant];
    QUEUEOCCUPATION(QFlow) = QUEUEOCCUPATION(QFlow) - QSize;
    QUEUE{QFlow}(1,:) = [];
end
\end{lstlisting}

O campo \texttt{PacketFlow} é necessário
porque a fila~1 contém pacotes dos dois fluxos, ou seja, ao contrário de em \texttt{Sim2FlowsPrio} o Id do fluxo e o Id da queue não correspondem. A entrada mais antiga
da fila, indicada pelo índice~1, continua a ser retirada primeiro. Esta
sequência implementa FIFO dentro da fila partilhada.

O valor \texttt{prio = 1} indica que o fluxo~3 deve ter prioridade máxima.
Esta foi a configuração usada nas experiências seguintes.

\newpage

\section{Tarefa 2.b -- Resultados de simulação para o caso C}

Foram consideradas as taxas do caso~C,
$\lambda=[800,800,600]$~pps, a capacidade $C=10$~Mbps e o critério de
paragem $P=10^5$. O enunciado especifica
$f_1=f_2=1\,000\,000$~B. Portanto, foi usada a configuração
$f=[1\,000\,000,1\,000\,000]$~B.

Foram feitas 50 simulações e foram calculados intervalos de confiança de 90\%
para a percentagem de perdas (PL), o atraso médio (APD), o atraso máximo (MPD)
e o throughput transmitido (TT).

\begin{table}[htbp]
    \centering
    \caption{Resultados de simulação para a Tarefa 2.b (50 réplicas, IC de 90\%).}
    \label{tab:2b}
    \begin{tabular}{lccc}
        \toprule
        Parâmetro & Fluxo 1 & Fluxo 2 & Fluxo 3 \\
        \midrule
        PL (\%)   & $0 \pm 0$ & $0 \pm 0$ & $0 \pm 0$ \\
        APD (ms)  & $10.9033 \pm 0.22786$ &
                     $10.3582 \pm 0.220268$ &
                     $1.10208 \pm 0.001228$ \\
        MPD (ms)  & $72.053 \pm 3.27981$ &
                     $71.8143 \pm 3.30082$ &
                     $4.58627 \pm 0.097721$ \\
        TT (Mbps) & $5.1154 \pm 0.006348$ &
                     $0.79995 \pm 0.000898$ &
                     $3.59802 \pm 0.005655$ \\
        \bottomrule
    \end{tabular}
\end{table}

\subsection*{Análise dos resultados}

Não houve perdas em nenhum dos três fluxos. Isto mostra que, neste cenário, as
filas não chegaram a atingir o limite definido.

O efeito da prioridade nota-se principalmente no atraso do fluxo~3. O seu APD é
$1.10208$~ms, enquanto os fluxos~1 e~2 apresentam $10.9033$~ms e
$10.3582$~ms, respetivamente. Quando chega um pacote do fluxo~3, ele passa à
frente dos pacotes dos fluxos~1 e~2 que ainda estão na fila. Por isso, o
fluxo~3 espera sobretudo pelo pacote que já estava a ser transmitido.

Na Tarefa~1.b todos os fluxos eram atendidos pela ordem de chegada e o APD do
fluxo~3 era aproximadamente $7.296$~ms. Com prioridades, esse valor desceu
para $1.102$~ms. Em contrapartida, os fluxos~1 e~2 passaram a esperar mais,
com APD de cerca de $10.9$~ms e $10.4$~ms. A prioridade melhora um fluxo, mas
faz aumentar o tempo de espera dos outros.

Os valores de TT permanecem praticamente iguais aos da Tarefa~1.b:
aproximadamente $5.12$, $0.80$ e $3.60$~Mbps. Isto é esperado, porque não
existem perdas e a prioridade apenas altera a ordem de transmissão.

O mesmo acontece com o MPD: o valor do fluxo~3 desce para cerca de $4.59$~ms,
enquanto os valores dos fluxos~1 e~2 ficam próximos de $72$~ms.

\newpage

\section{Tarefa 2.c -- Modelo M/G/1 com prioridades}

Para comparar os resultados da simulação com um modelo teórico, foi usado um
modelo M/G/1 com prioridades. Para cada fluxo calculou-se o tempo médio de
serviço $E[S_i]$, o segundo momento $E[S_i^2]$ e a utilização
$\rho_i=\lambda_iE[S_i]$.

Primeiro calcula-se o tempo de espera residual:
\[
W_0=\frac{1}{2}\sum_i\lambda_iE[S_i^2],
\]
Para o fluxo de maior prioridade, o tempo médio de espera é
\[
W_{Q,\mathrm{alta}}=\frac{W_0}{1-\rho_{\mathrm{alta}}},
\]
e para os fluxos de menor prioridade é
\[
W_{Q,\mathrm{baixa}}=
\frac{W_0}{(1-\rho_{\mathrm{alta}})(1-\rho_{\mathrm{total}})}.
\]

Neste caso, o fluxo 3 é o fluxo de maior prioridade e os fluxos~1 e~2 têm
menor prioridade. O atraso total de cada fluxo é a soma do tempo de espera com
o seu tempo médio de transmissão:
$D_i=W_{Q,i}+E[S_i]$.

\begin{lstlisting}[basicstyle=\ttfamily\scriptsize]
rho = y .* meanService;
totalRho = sum(rho);

rhoHigh = rho(3);
W0 = 0.5 * sum(y .* secondMomentService);

queueingDelayHigh = W0 / (1 - rhoHigh);
queueingDelayLow = ...
    W0 / ((1 - rhoHigh) * (1 - totalRho));

averagePacketDelayMs = 1000 * ...
    [queueingDelayLow + meanService(1), ...
     queueingDelayLow + meanService(2), ...
     queueingDelayHigh + meanService(3)];
\end{lstlisting}

Os resultados teóricos foram:

\begin{table}[htbp]
    \centering
    \caption{Resultados teóricos do modelo M/G/1 com prioridades.}
    \label{tab:2c-theory}
    \begin{tabular}{lcc}
        \toprule
        Fluxo & APD teórico (ms) & TT teórico (Mbps) \\
        \midrule
        1 & 11.114757 & 5.120000 \\
        2 & 10.574757 & 0.800000 \\
        3 & 1.102788  & 3.600000 \\
        \bottomrule
    \end{tabular}
\end{table}

\subsection*{Comparação com a simulação}

\begin{table}[htbp]
    \centering
    \caption{Comparação entre o modelo teórico e a simulação da Tarefa 2.b.}
    \label{tab:2c-comparison}
    \begin{tabular}{lccc}
        \toprule
        Fluxo & APD teórico (ms) & APD simulado (ms) & Erro relativo \\
        \midrule
        1 & 11.114757 & $10.9033 \pm 0.22786$ & 1.94\% \\
        2 & 10.574757 & $10.3582 \pm 0.220268$ & 2.09\% \\
        3 & 1.102788  & $1.10208 \pm 0.001228$ & 0.06\% \\
        \bottomrule
    \end{tabular}
\end{table}

Os valores teóricos e simulados são próximos. A maior diferença aparece nos
fluxos~1 e~2, porque o seu atraso depende tanto da chegada de pacotes do
fluxo~3 como da chegada dos pacotes do outro fluxo de menor prioridade. No
fluxo~3, os dois valores são praticamente iguais. As pequenas diferenças
resultam da aleatoriedade das simulações.

\section{Tarefa 2.d -- Atraso do fluxo 3 nos casos A, B e C}

Tal como na Tarefa~1.d, foram simulados os casos A, B e C, com 50 réplicas,
$P=10^5$ e $f_1=f_2=1$~MB. Nesta tarefa, o fluxo~3 tem prioridade máxima.

\begin{table}[htbp]
    \centering
    \caption{Atraso do fluxo 3 com prioridade (50 réplicas, IC de 90\%).}
    \label{tab:2d}
    \begin{tabular}{lcc}
        \toprule
        Caso & APD (ms) & MPD (ms) \\
        \midrule
        A & $0.938691 \pm 0.001047$ & $4.625990 \pm 0.115557$ \\
        B & $1.022540 \pm 0.001097$ & $4.529310 \pm 0.107777$ \\
        C & $1.103810 \pm 0.001472$ & $4.451330 \pm 0.099898$ \\
        \bottomrule
    \end{tabular}
\end{table}


\newpage

\begin{figure}[htbp]
    \centering
    \IfFileExists{Imagens/APD2.png}%
      {\includegraphics[width=0.75\linewidth]{Imagens/APD2.png}}%
      {\fbox{\parbox[c][3.5cm][c]{0.7\linewidth}{\centering\small
        Atraso médio do fluxo~3 com prioridade (adicionar
        \texttt{Imagens/APD2.png})}}}
    \caption{Atraso médio do fluxo 3 com prioridade nos casos A, B e C.}
    \label{fig:2d-apd}
\end{figure}

\begin{figure}[htbp]
    \centering
    \IfFileExists{Imagens/MPD2.png}%
      {\includegraphics[width=0.75\linewidth]{Imagens/MPD2.png}}%
      {\fbox{\parbox[c][3.5cm][c]{0.7\linewidth}{\centering\small
        Atraso máximo do fluxo~3 com prioridade (adicionar
        \texttt{Imagens/MPD2.png})}}}
    \caption{Atraso máximo do fluxo 3 com prioridade nos casos A, B e C.}
    \label{fig:2d-mpd}
\end{figure}

\subsection*{Justificação e comparação com a Tarefa 1.d}

O APD do fluxo~3 aumenta com a carga, de $0.939$~ms no caso~A para
$1.104$~ms no caso~C. Mesmo assim, este valor continua muito abaixo do
resultado da Tarefa~1.d, em que o APD no caso~C era aproximadamente
$7.296$~ms.

Isto acontece porque, sem prioridades, o fluxo~3 tinha de esperar por todos os
pacotes que estavam à sua frente. Com a nova regra, depois de acabar o pacote
que está a ser transmitido, o próximo pacote do fluxo~3 é escolhido antes dos
pacotes dos fluxos~1 e~2.

Os valores de MPD ficam perto de $4.5$~ms nos três casos. Na Tarefa~1.d, o MPD
do fluxo~3 chegava a aproximadamente $44.7$~ms no caso~C. Portanto, a
prioridade também reduz bastante os piores atrasos do fluxo~3.

\subsection*{Conclusões}

\begin{itemize}
    \item A prioridade reduz bastante o APD e o MPD do fluxo~3.
    \item O efeito é mais importante no caso~C, que tem maior carga.
    \item A redução do atraso do fluxo~3 é conseguida aumentando o tempo de
    espera dos fluxos~1 e~2.
\end{itemize}


\newpage

\section{Tarefa 3.a -- Simulador Deficit Round Robin para três fluxos}

Na Tarefa~3 foi desenvolvida a função \texttt{Sim3FlowsDrrT}, a partir do
simulador DRR com dois fluxos do Guião~2. A versão desenvolvida mantém uma fila
por fluxo e uma fila final DRR comum aos três fluxos.

\subsection*{Alterações introduzidas}

A adaptação dos contadores, da inicialização dos eventos e da geração dos
tamanhos dos pacotes segue o mesmo procedimento descrito na Tarefa~1 para a
introdução do terceiro fluxo. A alteração própria da Tarefa~3 está na função
\texttt{enterDRR}: a
ordem percorrida pela ronda passa a incluir o fluxo~3.

Tal como em Sim2FlowsDrr, em cada ronda, o algoritmo adiciona o quantum ao crédito do fluxo considerado
e move para a fila final tantos pacotes consecutivos quanto os que cabem no
crédito atual. O pacote que já estiver a ser transmitido não é interrompido.
Não houve necessidade de muitas alterações, uma vez que o código suporta
facilmente a adição de novas filas.

\begin{lstlisting}[basicstyle=\ttfamily\scriptsize]
if inv
    flowOrder = [3 2 1];
else
    flowOrder = [1 2 3];
end

for Flow = flowOrder
    if QUEUEOCCUPATION(Flow) == 0
        credit(Flow) = 0;
    else
        credit(Flow) = credit(Flow) + quantum(Flow);
        while QUEUEOCCUPATION(Flow) > 0 && QUEUE{Flow}(1,1) <= credit(Flow)
            PacketSize = QUEUE{Flow}(1,1);
            FINALQUEUE = [FINALQUEUE; QUEUE{Flow}(1,:)];
            QUEUE{Flow}(1,:) = [];
            QUEUEOCCUPATION(Flow) = QUEUEOCCUPATION(Flow) - PacketSize;
            credit(Flow) = credit(Flow) - PacketSize;
        end
    end
end
\end{lstlisting}
\newpage
\section{Tarefa 3.b -- DRR no caso C}

No caso~C consideraram-se
$\lambda=[800,800,600]$~pps, $C=10$~Mbps,
$f_1=f_2=f_3=1\,000\,000$~B e quantum igual a $1500$~B para cada fluxo.
Cada experiência foi executada 50 vezes, com $P=10^5$ e IC de 90\%.

\begin{table}[htbp]
    \centering
    \caption{Resultados DRR no caso C com quantum igual a 1500~B por fluxo.}
    \label{tab:3b}
    \begin{tabular}{lccc}
        \toprule
        Parâmetro & Fluxo 1 & Fluxo 2 & Fluxo 3 \\
        \midrule
        PL (\%) & $0 \pm 0$ & $0 \pm 0$ & $0 \pm 0$ \\
        APD (ms) & $10.0605 \pm 0.26231$ &
                   $2.09316 \pm 0.007437$ &
                   $4.02993 \pm 0.028119$ \\
        MPD (ms) & $71.4056 \pm 3.47416$ &
                   $6.81310 \pm 0.13718$ &
                   $24.2956 \pm 0.707198$ \\
        TT (Mbps) & $5.11136 \pm 0.006667$ &
                    $0.799958 \pm 0.000967$ &
                    $3.60206 \pm 0.004944$ \\
        \bottomrule
    \end{tabular}
\end{table}

Não se verificaram perdas no caso~C. O fluxo~2 apresenta o menor atraso médio,
pois os seus pacotes têm apenas 100--150~B. O fluxo~3 apresenta um atraso
intermédio e o fluxo~1 o maior atraso, consequência não apenas do
DRR, mas também da maior variação dos tamanhos dos seus pacotes, que podem
atingir 1500~B.

O throughput mantém-se próximo dos valores oferecidos pelos fluxos:
$5.11$, $0.80$ e $3.60$~Mbps. O valor muito próximo da carga teórica também
indica que a distribuição DRR não perde capacidade útil neste cenário.

\section{Tarefa 3.c -- DRR no caso congestionado D}

Mantiveram-se os 50 resultados por réplica, $P=10^5$ e filas de $1$~MB.
No caso~D considerou-se
$\lambda=[1000,1000,1000]$~pps e o quantum igual a $1500$~B para os três
fluxos.

\begin{table}[htbp]
    \centering
    \caption{Resultados DRR no caso D com quantum igual a 1500~B por fluxo.}
    \label{tab:3c}
    \begin{tabular}{lccc}
        \toprule
        Parâmetro & Fluxo 1 & Fluxo 2 & Fluxo 3 \\
        \midrule
        PL (\%) & $21.2683 \pm 0.079907$ &
                  $0 \pm 0$ &
                  $20.8901 \pm 0.091994$ \\
        APD (ms) & $1651.96 \pm 1.18139$ &
                   $2.78703 \pm 0.001160$ &
                   $1618.63 \pm 1.98380$ \\
        MPD (ms) & $1790.80 \pm 0.580693$ &
                   $5.50001 \pm 0.033551$ &
                   $1792.41 \pm 0.578203$ \\
        TT (Mbps) & $4.50011 \pm 0.000702$ &
                    $0.999796 \pm 0.001337$ &
                    $4.49990 \pm 0.000655$ \\
        \bottomrule
    \end{tabular}
\end{table}

Neste cenário, a carga oferecida excede a capacidade da ligação. Os fluxos~1
e~3 sofrem perdas semelhantes, enquanto o fluxo~2 não perde pacotes. Isto
acontece porque os pacotes do fluxo~2 são muito pequenos e cabem facilmente
no seu quantum de 1500~B. Em consequência, o seu atraso é muito reduzido.

Os fluxos~1 e~3 apresentam atrasos muito elevados, da ordem de 1600--1650~ms,
devido à acumulação nas respetivas filas após se atingir o limite de
transmissão. Os resultados do TT mostram que a ligação distribui aproximadamente
4.5~Mbps ao fluxo~1, 1~Mbps ao fluxo~2 e 4.5~Mbps ao fluxo~3.

\section{Tarefa 3.d -- DRR no caso D com quantum assimétrico}

Mantiveram-se as condições do caso~D, mas o quantum foi alterado para
$q=[1500,1500,4500]$~B. Assim, o fluxo~3 tem uma parcela três vezes superior
à dos fluxos~1 e~2 por ronda.

\begin{table}[htbp]
    \centering
    \caption{Resultados DRR no caso D com valores de quantum de 1500, 1500 e 4500~B.}
    \label{tab:3d}
    \begin{tabular}{lccc}
        \toprule
        Parâmetro & Fluxo 1 & Fluxo 2 & Fluxo 3 \\
        \midrule
        PL (\%) & $42.4007 \pm 0.117800$ &
                  $0 \pm 0$ &
                  $0 \pm 0$ \\
        APD (ms) & $2520.09 \pm 6.28820$ &
                   $3.87587 \pm 0.006908$ &
                   $6.36000 \pm 0.028581$ \\
        MPD (ms) & $2942.72 \pm 14.8186$ &
                   $14.5760 \pm 0.257080$ &
                   $27.5605 \pm 0.901622$ \\
        TT (Mbps) & $2.99691 \pm 0.007330$ &
                    $0.999819 \pm 0.001069$ &
                    $6.00290 \pm 0.007219$ \\
        \bottomrule
    \end{tabular}
\end{table}

O aumento do quantum do fluxo~3 reduz o seu atraso e elimina as perdas deste
fluxo. O TT passa a ser aproximadamente 6~Mbps no fluxo~3, enquanto o fluxo~2
continua a transmitir cerca de 1~Mbps. O throughput restante atribuído ao
fluxo~1 é aproximadamente 3~Mbps.

Esta alteração prejudica claramente o fluxo~1: a percentagem de perdas aumenta
de aproximadamente 21.27\% na Tarefa~3.c para 42.40\% na Tarefa~3.d, e o APD
aumenta de cerca de 1652~ms para 2520~ms. Portanto, o peso adicional do
fluxo~3 é atribuído às custas da parcela efetiva disponível ao fluxo~1.

\section{Tarefa 3.e -- Throughput ideal por weighted max-min fairness}

Os valores máximos de throughput oferecidos pelos três fluxos no caso~D são:

\[
R_1=1000\times800\times\frac{8}{10^6}=6.4~\text{Mbps},
\]
\[
R_2=1000\times125\times\frac{8}{10^6}=1.0~\text{Mbps},
\]
\[
R_3=1000\times750\times\frac{8}{10^6}=6.0~\text{Mbps}.
\]

A soma dos valores máximos de throughput é $13.4$~Mbps, superior à capacidade de 10~Mbps,
logo o sistema é congestionado.

\subsection*{Valores de quantum iguais}

Para valores de quantum iguais, os fluxos~1 e~3 têm igual peso e o fluxo~2 não requer
mais de 1~Mbps. Como o seu máximo é 1~Mbps, recebe primeiro 1~Mbps. Os restantes
9~Mbps são divididos igualmente pelos fluxos~1 e~3:

\[
T_{\text{ideal}}=[4.5,\,1.0,\,4.5]\ \text{Mbps}.
\]

Os resultados simulados são
$[4.50011,\,0.999796,\,4.49990]$~Mbps, praticamente iguais aos ideais.

\subsection*{Valores de quantum proporcionais a 1:1:3}

Com valores de quantum $[1500,1500,4500]$, os pesos relativos são $[1,1,3]$. A
distribuição ideal seria $[2,2,6]$~Mbps, mas o fluxo~2 tem máximo de 1~Mbps e
o fluxo~3 tem máximo de 6~Mbps.

Após atribuir a todos os fluxos a sua capacidade limitada:
\[
10-1-6=3~\text{Mbps},
\]
esta parcela restante é atribuída ao fluxo~1. Assim:

\[
T_{\text{ideal}}=[3.0,\,1.0,\,6.0]\ \text{Mbps}.
\]

Os resultados simulados são
$[2.99691,\,0.999819,\,6.00290]$~Mbps, confirmando uma forte concordância com
os valores teóricos.

\subsection*{Conclusões da Tarefa 3}

\begin{itemize}
    \item No caso~C, o sistema não apresenta perdas e o DRR transmite
    aproximadamente o mesmo throughput que em FIFO ou Priority Queuing.
    \item No caso~D, com valores de quantum iguais, o sistema distribui aproximadamente
    $[4.5,1,4.5]$~Mbps.
    \item No caso~D com o quantum $4500$~B no fluxo~3, a divisão passa a ser
    aproximadamente $[3,1,6]$~Mbps.
    \item A alteração dos valores de quantum permite controlar a parcela de serviço dos
    fluxos. O benefício do fluxo~3 é obtido sobretudo à custa do fluxo~1.
    \item Os resultados simulados aproximam-se dos valores ideais de
    weighted max-min fairness, com diferenças inferiores a 0.1\% nos três
    fluxos apresentados.
\end{itemize}

\end{document}


